\documentclass[uplatex,a4paper,12pt]{jsarticle}
\usepackage{amsmath,amssymb}
\usepackage{enumitem}

\begin{document}

\title{線形代数 確認テスト}
\author{}
\date{}
\maketitle

\section*{注意}
\begin{itemize}
\item 全10問，各10点，合計100点とする．
\item 60点以上を合格とする．
\item 計算問題では，必要に応じて途中式を書くこと．
\item 試験時間は15分とする．
\end{itemize}

\section*{問題}

\begin{enumerate}

\item 次のうち，ベクトルを表しているものを1つ選べ．
\begin{enumerate}[label=(\alph*)]
\item $5$
\item $(1,2,3)^T$
\item $\begin{pmatrix}1&2\\3&4\end{pmatrix}$
\end{enumerate}

\item ベクトル
\[
\mathbf{x}=(3,4)^T
\]
のユークリッドノルム $\|\mathbf{x}\|_2$ を求めよ．

\item ベクトル
\[
\mathbf{a}=(1,2,3)^T,\quad
\mathbf{b}=(4,0,2)^T
\]
の内積 $\mathbf{a}\cdot\mathbf{b}$ を求めよ．

\item 次の行列について，$A+B$ を求めよ．
\[
A=
\begin{pmatrix}
1&2\\
3&4
\end{pmatrix},
\quad
B=
\begin{pmatrix}
5&6\\
7&8
\end{pmatrix}
\]

\item 次の行列の転置行列 $A^T$ を求めよ．
\[
A=
\begin{pmatrix}
1&2&3\\
4&5&6
\end{pmatrix}
\]

\item 次の行列積 $AB$ を求めよ．
\[
A=
\begin{pmatrix}
1&2\\
3&4
\end{pmatrix},
\quad
B=
\begin{pmatrix}
2&0\\
1&3
\end{pmatrix}
\]

\item 次の行列について，アダマール積 $A\circ B$ を求めよ．
\[
A=
\begin{pmatrix}
1&2\\
3&4
\end{pmatrix},
\quad
B=
\begin{pmatrix}
2&0\\
1&3
\end{pmatrix}
\]

\item 次のうち，行列積について常に成り立つものをすべて選べ．
\begin{enumerate}[label=(\alph*)]
\item $AB=BA$
\item $A(BC)=(AB)C$
\item $A(B+C)=AB+AC$
\end{enumerate}

\item 次の行列のランクを求めよ．
\[
A=
\begin{pmatrix}
1&2\\
2&4
\end{pmatrix}
\]

\item 次のうち，固有値・固有ベクトルの関係式として正しいものを1つ選べ．
\begin{enumerate}[label=(\alph*)]
\item $Ap=\lambda p$
\item $A+p=\lambda$
\item $A=\lambda+p$
\end{enumerate}

\end{enumerate}

\end{document}